\chapter{Decoder-only Transformer의 정확한 순전파}
\label{chap:phase2}

\section{결론부터: 다음 토큰 분포는 합성함수의 출력이다}

decoder-only Transformer의 한 번의 순전파는 token sequence를 vocabulary 전체의 logits로 변환하는 결정론적 함수다.
dropout을 끄고 가중치와 수치 연산을 고정하면, 입력 token IDs가 같을 때 logits도 같다.
다음 token을 실제로 고르는 sampling 단계에는 별도의 난수가 들어갈 수 있다.
\index{autoregressive generation@자기회귀 생성(autoregressive generation)}

pre-normalization을 사용하는 순차형 block은 attention update와 MLP update를 차례로 residual stream에 더한다.
\begin{align}
  \mat{U}^{(\ell)}
    &=\mat{X}^{(\ell)}+\operatorname{Attn}^{(\ell)}
       \left(\operatorname{Norm}^{(\ell)}_1(\mat{X}^{(\ell)})\right),
       \label{eq:attention-residual-update} \\
  \mat{X}^{(\ell+1)}
    &=\mat{U}^{(\ell)}+\operatorname{MLP}^{(\ell)}
       \left(\operatorname{Norm}^{(\ell)}_2(\mat{U}^{(\ell)})\right).
  \label{eq:prenorm-block}
\end{align}
마지막 layer 뒤에서는
\begin{align}
  \mat{H}&=\operatorname{Norm}_{f}(\mat{X}^{(L)}),\\
  \mat{Z}&=\mat{H}\mat{W}_{U}^{\mathsf T}+\vect{b}_{U},\\
  p(x_{T+1}=v\mid x_{1:T})&=\softmax(\mat{Z}_{:,T,:})_v
\end{align}
를 계산한다.
\Cref{eq:attention-residual-update}은 현재 residual stream $\mat{X}^{(\ell)}$을 normalize하여 multi-head attention에 넣고, attention이 계산한 update를 원래의 residual stream에 더하는 식이다.
실제 구현은 bias를 생략하거나 embedding과 unembedding weight를 공유할 수 있다.
원래 Transformer의 모든 변형이 \cref{eq:prenorm-block}과 같은 배치를 사용하는 것은 아니다 \citep{vaswani2017attention,xiong2020layernorm}.

\section{기호와 tensor shape}

순전파를 검산하려면 모든 수식에서 같은 축과 차원 기호를 사용해야 한다.
이 장에서는 다음 기호를 고정한다.

\begin{table}[htbp]
\centering
\caption{기본 차원과 의미}
\label{tab:forward-dimensions}
\begin{tabularx}{\textwidth}{@{}lYl@{}}
\toprule
기호 & 의미 & 대표 shape 또는 범위 \\
\midrule
$B$ & batch에 포함된 sequence 수 & 양의 정수 \\
$T$ & 현재 sequence 길이 & 양의 정수 \\
$V$ & vocabulary 크기 & 양의 정수 \\
$d$ & residual stream 차원, $d_{\text{model}}$ & 양의 정수 \\
$H$ & query 및 MHA의 head 수 & 양의 정수 \\
$d_h$ & 한 head의 차원 & 보통 $d/H$ \\
$d_{\mathrm{ff}}$ & MLP 중간 차원 & 모델별로 다름 \\
$L$ & Transformer block 수 & 양의 정수 \\
$\mat{X}^{(\ell)}$ & layer $\ell$ 입구의 residual stream & $\shape{B,T,d}$ \\
$\mat{Z}$ & 각 위치와 token의 logit & $\shape{B,T,V}$ \\
\bottomrule
\end{tabularx}
\end{table}

수식에서는 읽기 쉽도록 batch 축을 생략할 때가 있다.
그러나 실제 구현에서 batch 축이 사라지는 것은 아니다.
position $i$의 query는 destination, position $j$의 key와 value는 source라고 부르겠다.
이 이름을 고정하면 attention matrix의 행과 열을 뒤바꾸는 실수를 줄일 수 있다.

\section{토큰 ID에서 residual stream으로}

tokenizer가 문자열을 정수열 $(t_1,\ldots,t_T)$로 바꾸면 $t_i\in\{0,\ldots,V-1\}$이다.
여기서 $t_i$는 하나의 token ID이고 $V$는 vocabulary size다.
embedding matrix $\mat{E}\in\R^{V\times d}$에서 각 row를 선택하여
\begin{equation}
  \mat{X}^{(0)}_{b,i,:}=\mat{E}_{t_{b,i},:}
\end{equation}
를 얻는다.
여기서 $V$는 row의 수인 vocabulary size이고 $d$는 column의 수인 embedding dimension이다.
즉, 각 token ID는 $d$차원 learned vector 하나에 대응한다.
absolute positional embedding을 사용하는 모델은 position vector를 여기서 더할 수 있다.
RoPE를 사용하는 모델은 일반적으로 input embedding에 position vector를 더하지 않고 attention 안의 query와 key를 회전한다.
token ID를 embedding matrix의 row index로 사용하여 그 token ID에 대응하는 vector를 가져온다.
absolute positional embedding을 사용하면 첫 block 전에 position vector를 더하고, RoPE를 사용하면 각 attention layer에서 query와 key에 position-dependent rotation을 적용한다.
자세한 내용은 \cref{chap:phase3}에서 다룬다.

padding을 포함하는 batch에서는 padding 위치를 막는 mask가 causal mask와 별도로 필요할 수 있다.
반면 단일 sequence만 처리하며 padding이 없다면 causal mask만 있으면 된다.
tokenizer가 만든 token 경계와 모델이 처리하는 position은 문자나 단어 경계와 일치할 필요가 없다.

\begin{factbox}
한 position의 초기 벡터는 그 token ID에 대응하는 embedding이다.
첫 attention layer를 통과한 뒤에는 앞선 여러 position의 value가 섞일 수 있으므로, 같은 token ID도 문맥에 따라 서로 다른 residual vector를 갖는다.
\end{factbox}

\section{한 attention head의 모든 계산}

attention head는 각 destination position에서 causal mask가 허용하는 현재 또는 과거 source position 가운데 어디를 참고할지 정하고, 그 위치들의 value 정보를 섞어 가져오는 계산 단위다.
한 attention head의 계산은 다음 네 단계로 요약할 수 있다.

\begin{enumerate}
  \item query projection으로 destination position이 찾는 정보의 routing feature를 계산한다.
  \item key projection으로 각 source position이 제공할 수 있는 정보의 routing feature를 계산한다.
  \Needspace{3\baselineskip}
  \item query와 key를 비교하고 causal mask와 softmax를 적용하여 attention weights를 계산한다.
  \item attention weights로 source별 value를 가중합하여 head output을 만든다.
\end{enumerate}

이 설명은 Q, K와 V의 계산상 역할을 나타내는 intuition이며, 개별 vector가 반드시 사람이 이름 붙일 수 있는 의미 하나와 대응한다는 뜻은 아니다.
Q는 ``현재 destination이 어떤 정보를 찾는가'', K는 ``각 source가 어떤 정보를 제공할 수 있는가'', V는 ``선택된 source에서 실제로 전달할 내용이 무엇인가''라는 식으로 이해할 수 있다.

한 attention head는 정규화한 residual stream에서 query, key와 value를 각각 투영한다.
$\mat{R}=\operatorname{Norm}(\mat{X})\in\R^{B\times T\times d}$라고 두면 head $h$의 세 projection은 다음과 같다.
\begin{align}
  \mat{Q}^{h}&=\mat{R}\mat{W}^{h}_{Q}, &
  \mat{W}^{h}_{Q}&\in\R^{d\times d_h},\\
  \mat{K}^{h}&=\mat{R}\mat{W}^{h}_{K}, &
  \mat{W}^{h}_{K}&\in\R^{d\times d_h},\\
  \mat{V}^{h}&=\mat{R}\mat{W}^{h}_{V}, &
  \mat{W}^{h}_{V}&\in\R^{d\times d_h}.
\end{align}
따라서 $\mat{Q}^{h},\mat{K}^{h},\mat{V}^{h}$의 shape는 모두 $\shape{B,T,d_h}$다.

계산 순서를 수식과 연결하면 다음과 같다.

\begin{enumerate}
  \item residual stream $\mat{X}\in\R^{B\times T\times d}$를 normalize하여 $\mat{R}=\operatorname{Norm}(\mat{X})\in\R^{B\times T\times d}$를 만든다.
  \item 각 head는 같은 $\mat{R}$에 서로 다른 weight matrix를 곱하여 $\mat{Q}^{h}=\mat{R}\mat{W}^{h}_{Q}$, $\mat{K}^{h}=\mat{R}\mat{W}^{h}_{K}$와 $\mat{V}^{h}=\mat{R}\mat{W}^{h}_{V}$를 만든다.
\end{enumerate}

예를 들어 ``The cat sat on the mat because it was tired.''라는 문장에서 coreference를 수행하는 hypothetical head를 생각해 보자.
설명을 단순화하기 위해 ``cat''과 ``it''이 각각 하나의 token position에 대응한다고 가정한다.
``it''의 query는 ``내가 가리키는 대상이 누구인가''와 관련된 routing feature를 가질 수 있다.
``cat''의 key는 ``대명사가 참조할 수 있는 동물 또는 명사''와 관련된 routing feature를 가질 수 있다.
``it'' position의 query와 ``cat'' position의 key가 다른 허용된 source들의 key보다 잘 맞으면 softmax 뒤에서 ``cat'' position에 큰 attention weight가 배정될 수 있다.
이때 head가 가져오는 내용은 ``cat''의 raw residual vector가 아니라 해당 position의 value projection이다.
실제 tokenizer에서는 한 단어가 여러 token으로 분리될 수 있고, 실제 head가 반드시 이런 human-interpretable 역할을 수행한다는 보장도 없다.

destination $i$와 source $j$ 사이의 attention score는
\begin{equation}
  S^h_{b,i,j}=
  \frac{\langle\vect{q}^{h}_{b,i},\vect{k}^{h}_{b,j}\rangle}{\sqrt{d_h}}
  +M_{i,j}
\end{equation}
이다.
여기서 destination position $i$의 query와 source position $j$의 key가 얼마나 잘 맞는지를 dot product로 계산한다.
한 raw score의 절댓값만으로 attention weight가 크다고 결론낼 수 없으며, 같은 destination row의 다른 source scores와 함께 mask와 softmax를 적용한 뒤 비교해야 한다.
query와 key의 각 성분이 비슷한 분산을 가지고 약하게 상관된다고 가정하면 dot product의 분산은 대략 $d_h$에 비례하므로, $\sqrt{d_h}$로 나누어 score의 scale을 안정화한다.

causal mask는 decoder-only Transformer가 미래 token을 직접 참고하지 못하게 하며,
\begin{equation}
  M_{i,j}=\begin{cases}
  0,&j\le i,\\
  -\infty,&j>i
  \end{cases}
\end{equation}
로 두며, 실제 유한정밀도 코드에서는 dtype에 맞는 매우 작은 값을 사용할 수도 있다.
source 축에 row-wise softmax를 적용하면
\begin{equation}
  A^h_{b,i,j}=\frac{\exp S^h_{b,i,j}}
  {\sum_{k=1}^{T}\exp S^h_{b,i,k}},
  \qquad \sum_j A^h_{b,i,j}=1
\end{equation}
을 얻는다.
각 destination row에서 attention weights는 nonnegative이고 합이 1이지만, 특정 source가 의미론적으로 정답일 확률로 calibration된 값은 아니다.
masked position의 weight는 이상적인 실수 연산에서 0이다.
\index{causal attention@인과적 어텐션(causal attention)}
\index{QKV@query--key--value (QKV)}

한 head의 출력은 source별 value의 attention-weighted sum이다.
\begin{equation}
  \vect{o}^{h}_{b,i}=\sum_{j=1}^{T}A^h_{b,i,j}\vect{v}^{h}_{b,j}.
  \label{eq:head-weighted-sum}
\end{equation}
$\mat{A}^{h}$의 shape는 $\shape{B,T,T}$이고 $\mat{O}^{h}$의 shape는 $\shape{B,T,d_h}$다.
causal mask 때문에 position $i$의 출력은 $j\le i$인 value만 직접 사용한다.
즉, residual vector 자체를 가져오는 것이 아니라 $\vect{v}^{h}_{b,j}=\mat{R}_{b,j,:}\mat{W}^{h}_{V}$라는 value projection을 거친 정보를 가져온다.

\begin{interpretbox}{``정보를 가져온다''는 직관}
\Cref{eq:head-weighted-sum} 때문에 attention을 source position에서 value 정보를 가져오는 연산이라고 설명할 수 있다.
다만 어떤 source를 선택하는지는 query와 key의 문맥 의존적 계산이며, 전달되는 내용도 원 residual vector 자체가 아니라 value projection이다.
attention weight만 보아서는 output projection 이후의 효과를 알 수 없다.
\end{interpretbox}

\section{multi-head attention과 output projection}

multi-head attention은 $H$개 head의 출력을 이어 붙인 뒤 output projection으로 residual 차원에 되돌린다.
즉, 여러 head의 결과를 하나로 합쳐서 다시 residual stream의 차원 $d$로 되돌린다.
head 출력을 마지막 축에서 이어 붙이면
\begin{equation}
  \mat{O}_{\mathrm{cat}}=\operatorname{Concat}
  (\mat{O}^{1},\ldots,\mat{O}^{H})
  \in\R^{B\times T\times Hd_h}.
\end{equation}
$Hd_h=d$인 일반적인 설정에서는 output projection $\mat{W}_{O}\in\R^{d\times d}$를 곱하여
\begin{equation}
  \Delta\mat{X}_{\mathrm{attn}}=
  \mat{O}_{\mathrm{cat}}\mat{W}_{O}
  \in\R^{B\times T\times d}
\end{equation}
를 얻는다.
이를 현재 residual stream에 더할 수 있는 shape로 되돌리는 과정이다.

$\mat{W}_{O}$를 head별 row block으로 나누면 전체 출력은 head별 residual update의 합으로도 쓸 수 있다.
\begin{equation}
  \Delta\mat{X}_{\mathrm{attn}}
  =\sum_{h=1}^{H}\mat{O}^{h}\mat{W}^{h}_{O}.
\end{equation}
따라서 concatenation 구현과 head별 합 표현은 서로 다른 모델이 아니라 같은 선형연산의 두 표기다 \citep{elhage2021framework}.

예를 들어, $d=4096$이고 head가 $H=32$라 하면 $d_h=4096/32=128$이다.
즉, 각 head는 token마다 128차원짜리 결과를 낸다.
이 32개 head output을 이어 붙이면
\begin{equation}
  \mat{O}_{\mathrm{cat}}
  =\operatorname{Concat}(\mat{O}^{1},\ldots,\mat{O}^{32}),
  \qquad
  \operatorname{shape}(\mat{O}_{\mathrm{cat}})
  =\shape{B,T,32\cdot128}
  =\shape{B,T,4096}
\end{equation}
가 된다.
$\mat{W}_{O}\in\R^{4096\times4096}$는 output projection matrix다.
이 $\mat{W}_{O}$를 $\mat{O}_{\mathrm{cat}}$에 곱하면
\begin{equation}
  \Delta\mat{X}_{\mathrm{attn}}
  =\mat{O}_{\mathrm{cat}}\mat{W}_{O},
  \qquad
  \operatorname{shape}(\Delta\mat{X}_{\mathrm{attn}})
  =\shape{B,T,4096}
\end{equation}
가 된다.
따라서 단순히 이어 붙인 32개 head output은 output projection에서 서로 섞이고 residual stream에 쓸 수 있는 4096차원 update로 변환된다.

attention pattern $\mat{A}^{h}$는 query와 key에서 계산되므로 input에 의존하고, dot product와 softmax를 거치므로 전체 attention은 input에 대한 단순한 선형변환이 아니다.
pre-normalization block에서는 $\mat{Q}^{h}=\mat{R}\mat{W}^{h}_{Q}$와 $\mat{K}^{h}=\mat{R}\mat{W}^{h}_{K}$이므로, $\mat{R}$이 달라지면 attention pattern도 달라질 수 있다.

head $h$의 QK와 OV weight product를
\begin{equation}
  \mat{W}^{h}_{QK}
  =\mat{W}^{h}_{Q}(\mat{W}^{h}_{K})^{\mathsf T}
  \in\R^{d\times d},
  \qquad
  \mat{W}^{h}_{OV}
  =\mat{W}^{h}_{V}\mat{W}^{h}_{O}
  \in\R^{d\times d}
\end{equation}
로 정의할 수 있다.
그러면 attention score는
\begin{equation}
  S^{h}_{b,i,j}
  =\frac{\mat{R}_{b,i,:}\mat{W}^{h}_{QK}\mat{R}_{b,j,:}^{\mathsf T}}
         {\sqrt{d_h}}
  +M_{i,j}
\end{equation}
로 쓸 수 있고, QK circuit은 destination과 source 사이의 routing score를 결정한다.
OV circuit은 선택된 source의 normalized residual vector를 residual stream에 쓸 방향으로 변환하며, head별 update는
\begin{equation}
  \Delta\mat{X}^{h}_{\mathrm{attn},b,i,:}
  =\sum_{j=1}^{T}
   A^{h}_{b,i,j}\mat{R}_{b,j,:}\mat{W}^{h}_{OV}
\end{equation}
로 쓸 수 있다.

특정 input에서 미리 계산한 attention pattern $\mat{A}^{h}$를 고정했다고 하자.
이때 head output은 $\mat{O}^{h}=\mat{A}^{h}\mat{V}^{h}=\mat{A}^{h}\mat{R}\mat{W}^{h}_{V}$다.
output projection까지 포함한 head별 update는 $\Delta\mat{X}^{h}_{\mathrm{attn}}=\mat{A}^{h}\mat{R}\mat{W}^{h}_{V}\mat{W}^{h}_{O}=\mat{A}^{h}\mat{R}\mat{W}^{h}_{OV}$다.
$\mat{A}^{h}$, $\mat{W}^{h}_{V}$와 $\mat{W}^{h}_{O}$를 모두 고정하면 이 경로는 normalized activation $\mat{R}$에 대해 선형이다.
그러나 $\mat{R}=\operatorname{Norm}(\mat{X})$이므로 normalization까지 포함한 raw residual stream $\mat{X}$에서 update로 가는 전체 함수가 선형이라는 뜻은 아니다.
mechanistic interpretability에서 QK circuit을 ``어디를 볼 것인가'', OV circuit을 ``선택한 위치에서 어떤 정보를 residual stream에 쓸 것인가''로 구분하는 이유가 여기에 있다.

\section{MLP update와 잔차 덧셈}

sequential pre-norm block에서 MLP는 attention update가 더해진 중간 residual state를 읽는다.
attention sublayer의 출력과 MLP가 읽는 normalized activation을 다음과 같이 두자.
\begin{align}
  \mat{U}
  &=\mat{X}+\Delta\mat{X}_{\mathrm{attn}},\\
  \mat{R}_{\mathrm{mlp}}
  &=\operatorname{Norm}_{2}(\mat{U})
\end{align}
각 position에 독립적으로 적용되는 기본 MLP는
\begin{align}
  \mat{P}
  &=\mat{R}_{\mathrm{mlp}}\mat{W}_{\mathrm{in}}+\vect{b}_{\mathrm{in}},\\
  \mat{G}
  &=\phi(\mat{P}),\\
  \Delta\mat{X}_{\mathrm{mlp}}
  &=\mat{G}\mat{W}_{\mathrm{out}}+\vect{b}_{\mathrm{out}}
\end{align}
로 쓸 수 있다.
여기서 $\mat{R}_{\mathrm{mlp}}\in\R^{B\times T\times d}$, $\mat{W}_{\mathrm{in}}\in\R^{d\times d_{\mathrm{ff}}}$, $\mat{W}_{\mathrm{out}}\in\R^{d_{\mathrm{ff}}\times d}$다.
$\mat{P}$와 $\mat{G}$의 shape는 $\shape{B,T,d_{\mathrm{ff}}}$이고 $\Delta\mat{X}_{\mathrm{mlp}}$의 shape는 $\shape{B,T,d}$다.
bias vector는 batch와 position 축으로 broadcast된다.

같은 MLP weight가 모든 position에 적용되지만, position $i$의 출력은 오직 $\mat{R}_{\mathrm{mlp},b,i,:}$에서 계산된다.
따라서 MLP 자체는 서로 다른 token position을 직접 섞지 않는다.
다만 $\mat{R}_{\mathrm{mlp},b,i,:}$는 앞선 attention과 이전 block들이 만든 contextualized vector이므로 다른 position에서 온 정보가 이미 들어 있을 수 있다.
이 장에서 다루는 표준 block에서는 position 사이의 직접적인 정보 전달은 attention에서 일어나고, MLP는 각 position 안에서 feature를 변환한다.
\index{MLP@MLP}

예를 들어 $d=4096$이고 $d_{\mathrm{ff}}=14{,}336$이라 하자.
$\mat{R}_{\mathrm{mlp}}$의 shape가 $\shape{B,T,4096}$이면 input projection은 각 token의 4096차원 vector를 14,336차원 pre-activation으로 보낸다.
따라서 $\mat{P}$와 $\mat{G}$의 shape는 $\shape{B,T,14336}$다.
down projection $\mat{W}_{\mathrm{out}}\in\R^{14336\times4096}$을 적용하면 $\Delta\mat{X}_{\mathrm{mlp}}$의 shape는 다시 $\shape{B,T,4096}$이 된다.
이렇게 residual stream에 더할 수 있도록 마지막 차원을 $d$로 되돌린다.

$\phi$가 identity이고 gating이 없다면 두 affine projection의 합성은 하나의 affine map으로 합칠 수 있다.
nonlinearity 또는 gating이 있으면 MLP는 입력에 따라 서로 다른 intermediate feature를 활성화하므로 하나의 고정된 affine map으로 축약할 수 없다.
현대 모델의 SwiGLU는 두 input projection을 element-wise gating으로 결합하므로 이 기본식과 parameterization이 다르며, 정확한 식은 \cref{chap:phase3}에서 다룬다.

한 sequential pre-norm block의 전체 residual update는
\begin{align}
  \mat{U}
  &=\mat{X}
    +\operatorname{Attn}\!\left(\operatorname{Norm}_{1}(\mat{X})\right),\\
  \mat{X}_{\mathrm{next}}
  &=\mat{U}
    +\operatorname{MLP}\!\left(\operatorname{Norm}_{2}(\mat{U})\right)
\end{align}
이다.
attention update가 먼저 $\mat{U}$에 반영되고, MLP는 이 갱신된 $\mat{U}$를 normalize한 값을 읽는다.
따라서 sequential block의 MLP update는 일반적으로 원래 $\mat{X}$만의 함수가 아니라 attention update에도 의존한다.
덧셈으로 인해 이전 residual state로 향하는 identity path가 항상 존재한다.
그러나 identity path가 있다는 사실만으로 이전 information이 이후 layer에서도 그대로 decode되거나 사용된다고 보장할 수는 없다.
다만 normalization과 이후 layer가 개별 update를 비선형적으로 다시 읽으므로, 마지막 logit에 대한 각 update의 총 효과를 단순히 독립적인 상수 기여로 볼 수는 없다.

\begin{factbox}
``hidden state가 layer마다 완전히 교체된다''는 설명보다 ``attention과 MLP가 residual stream에 update를 더한다''는 설명이 residual architecture의 계산을 더 직접적으로 반영한다.
그러나 identity path는 기존 information의 보존을 보장하지 않으며, normalization 위치와 parallel/sequential block 여부도 모델마다 다르므로 실제 코드를 확인해야 한다.
\end{factbox}

\section{마지막 정규화, unembedding과 logits}

마지막 block의 residual stream을 $\mat{X}^{(L)}\in\R^{B\times T\times d}$라고 하자.
final normalization과 unembedding 또는 LM head는
\begin{align}
  \mat{H}
  &=\operatorname{Norm}_{f}(\mat{X}^{(L)}),\\
  \mat{Z}
  &=\mat{H}\mat{W}_{U}^{\mathsf T}+\vect{b}_{U},
  &\mat{W}_{U}&\in\R^{V\times d}
\end{align}
를 계산한다.
$\mat{H}$의 shape는 $\shape{B,T,d}$이고 $\mat{Z}$의 shape는 $\shape{B,T,V}$다.
$\mat{W}_{U}$의 각 row는 residual dimension의 vector를 특정 vocabulary token의 logit으로 읽는 unembedding direction이다.
bias 사용 여부와 input embedding과의 weight tying 여부는 모델마다 다르다.

prompt 다음 token을 예측할 때는 각 sequence의 마지막 유효 position $T_b$에서 $\mat{Z}_{b,T_b,:}\in\R^{V}$를 선택한다.
이 vector에 softmax를 적용하면
\begin{equation}
  p(x_{T_b+1}=v\mid x_{1:T_b})
  =\softmax(\mat{Z}_{b,T_b,:})_v
\end{equation}
인 다음-token 조건부분포를 얻는다.

훈련에서는 causal mask 덕분에 여러 position의 logits를 한 번에 계산할 수 있다.
position $t$의 logit row는 같은 position의 input token $x_t$가 아니라 다음 token $x_{t+1}$을 target으로 사용한다.
padding이나 별도의 loss mask가 없다고 단순화하면 한 sequence의 negative log-likelihood는
\begin{equation}
  \mathcal{L}
  =-\sum_{t=1}^{T-1}
   \log\softmax(\mat{Z}_{t,:})_{x_{t+1}}
\end{equation}
로 쓸 수 있다.
실제 batch에서는 padding, special token과 학습 대상이 아닌 position을 loss mask로 제외할 수 있다.

position $t$에서 token $v$의 logit은 row-vector 관례로
\begin{equation}
  z_{t,v}
  =\langle\vect{h}_{t},\vect{w}_{U,v}\rangle+b_{U,v}
\end{equation}
이다.
residual update $\Delta\vect{h}$가 normalization을 거치지 않는 이상적인 선형 경로에서 주는 direct logit contribution은 $\langle\Delta\vect{h},\vect{w}_{U,v}\rangle$다.
실제 모델의 final normalization을 무시하면 오차가 생길 수 있으므로 direct logit attribution의 정의를 명시해야 한다.

예를 들어 $B=1$, $T=8$, $d=4096$, $V=128{,}256$이라 하자.
final normalization 결과 $\mat{H}$의 shape는 $\shape{1,8,4096}$이고 $\mat{W}_{U}^{\mathsf T}$의 shape는 $\shape{4096,128256}$이다.
두 tensor를 곱하면 모든 position의 logits $\mat{Z}$는 $\shape{1,8,128256}$이 된다.
마지막 유효 position을 선택하면 next-token logits의 shape는 $\shape{1,128256}$이다.

\begin{factbox}
logit은 softmax가 확률로 정규화하기 전의 score이며 확률이 아니다.
softmax는 vocabulary의 모든 logit을 함께 사용한다.
특정 token의 logit이 증가해도 다른 token의 logit이 더 많이 증가하면 그 token의 확률은 감소할 수 있다.
또한 final normalization은 residual vector의 모든 coordinate를 함께 읽으므로, 한 residual update의 정확한 logit 효과가 그 update와 unembedding direction의 내적만으로 항상 주어지는 것은 아니다.
\end{factbox}

\section{autoregressive generation}

한 번의 순전파가 분포 $p(x_{T+1}\mid x_{1:T})$를 만들면 decoder가 token $\hat{x}_{T+1}$을 선택한다.
대표적인 규칙은 다음과 같다.

\begin{itemize}
  \item greedy decoding은 $\argmax_v z_v$를 선택한다.
  \item temperature sampling은 $\tau>0$에 대해 $\softmax(\vect{z}/\tau)$에서 표본을 뽑는다.
  \item top-$k$ sampling은 logit이 큰 상위 $k$개 token만 남기고 나머지를 제외한 뒤 다시 정규화한다.
  \item nucleus 또는 top-$p$ sampling은 확률이 큰 순서로 token을 정렬하고 누적 확률이 기준 $p$ 이상이 되는 최소 후보 집합을 남긴 뒤 다시 정규화한다 \citep{holtzman2020degeneration}.
\end{itemize}

대표적인 sampling pipeline은 다음 순서로 이해할 수 있다.

\begin{enumerate}
  \item 현재 prefix를 모델에 넣어 마지막 유효 position의 logits를 계산한다.
  \item sampling을 사용하면 temperature로 logits의 상대적 scale을 조정한다.
  \item top-$k$ 또는 top-$p$를 사용하면 후보 집합 밖의 token을 제거한다.
  \item 남은 후보를 정규화한 분포에서 token을 뽑거나 greedy rule로 하나를 선택한다.
  \item 선택한 token을 prefix 뒤에 붙이고 다음 generation step을 시작한다.
\end{enumerate}

예를 들어 네 token $(A,B,C,D)$의 logits가 $\vect{z}=(3,2,1,0)$이라고 하자.
greedy decoding은 가장 큰 logit을 가진 $A$를 선택한다.
$\tau=1$인 원래 softmax 확률은 대략 $(0.644,0.237,0.087,0.032)$이므로 top-$2$는 $A$와 $B$만 후보로 남긴다.
top-$p$에서 $p=0.8$을 사용해도 $A$ 하나의 누적 확률은 0.8보다 작고 $A$와 $B$의 누적 확률은 약 0.881이므로 두 token을 남긴다.
$\tau=2$이면 확률은 대략 $(0.455,0.276,0.167,0.102)$로 평평해지고, $\tau=0.5$이면 대략 $(0.865,0.117,0.016,0.002)$로 가장 큰 logit에 더 집중된다.
sampling에서는 후보로 남은 token 가운데 반드시 가장 큰 logit의 token만 선택되는 것은 아니다.
temperature, top-$k$와 top-$p$는 해당 generation step의 model weight나 forward-pass 내부 activation을 바꾸는 것이 아니라 forward-pass가 만든 logits에서 decoding distribution을 구성하는 방법을 바꾼다.

선택한 token을 sequence 뒤에 붙여 같은 조건부 예측을 반복한다.
\begin{equation}
  p(x_{T+1:T+n}\mid x_{1:T})
  =\prod_{r=1}^{n}p(x_{T+r}\mid x_{1:T+r-1}).
\end{equation}
따라서 한 응답은 단 한 번의 ``전체 문장 출력''이 아니라 token별 조건부 선택의 연쇄다.
greedy에서는 선택이 결정론적이지만, 앞 단계에서 선택한 token이 이후 모든 조건부분포의 입력을 바꾸므로 초기 차이는 계속 전파될 수 있다.

구현에서 ``temperature 0''은 대개 0으로 나누어 softmax를 계산한다는 뜻이 아니라 sampling을 끄고 argmax를 선택한다는 설정 이름이다.
실제 구현은 prompt 전체를 처리하는 prefill과 새 token을 하나씩 처리하는 decode를 구분하고 KV cache로 과거 key와 value를 재사용할 수 있으며, 이 실행 경로는 \cref{chap:phase3}에서 다룬다.

\section{완전한 shape 추적 예제}

완전한 shape 추적은 token ID부터 next-token logits까지 각 축의 의미를 보존해야 한다.
\Cref{tab:shape-trace}는 $B=1$, $T=8$, $d=4096$, $H=32$, $d_h=128$인 MHA block을 가정하며, vocabulary와 MLP 크기는 예시로 각각 $128{,}256$, $14{,}336$을 사용한다.
이 표의 MLP는 위에서 정의한 하나의 input projection과 nonlinearity를 사용하는 기본 MLP이며, SwiGLU에서는 intermediate projection이 둘로 나뉜다.

\begin{table}[htbp]
\centering
\small
\caption{한 sequential pre-norm block과 LM head의 shape 추적}
\label{tab:shape-trace}
\begin{tabularx}{\textwidth}{@{}YYl@{}}
\toprule
대상 & 계산 & shape \\
\midrule
token IDs & tokenizer 출력 & $\shape{1,8}$ \\
embedding/residual $\mat{X}$ & $\mat{E}[\mat{t}]$ & $\shape{1,8,4096}$ \\
attention norm $\mat{R}_{\mathrm{attn}}$ & $\operatorname{Norm}_{1}(\mat{X})$ & $\shape{1,8,4096}$ \\
query & reshape$(\mat{R}_{\mathrm{attn}}\mat{W}_Q)$ & $\shape{1,32,8,128}$ \\
key & reshape$(\mat{R}_{\mathrm{attn}}\mat{W}_K)$ & $\shape{1,32,8,128}$ \\
value & reshape$(\mat{R}_{\mathrm{attn}}\mat{W}_V)$ & $\shape{1,32,8,128}$ \\
scores & $\mat{Q}\mat{K}^{\mathsf T}/\sqrt{128}+\mat{M}$ & $\shape{1,32,8,8}$ \\
weights & source 축 softmax & $\shape{1,32,8,8}$ \\
head output & $\mat{A}\mat{V}$ & $\shape{1,32,8,128}$ \\
concatenated output & transpose, reshape & $\shape{1,8,4096}$ \\
attention update & $\mat{O}_{\mathrm{cat}}\mat{W}_O$ & $\shape{1,8,4096}$ \\
attention residual $\mat{U}$ & $\mat{X}+\Delta\mat{X}_{\mathrm{attn}}$ & $\shape{1,8,4096}$ \\
MLP norm $\mat{R}_{\mathrm{mlp}}$ & $\operatorname{Norm}_{2}(\mat{U})$ & $\shape{1,8,4096}$ \\
MLP pre-activation $\mat{P}$ & $\mat{R}_{\mathrm{mlp}}\mat{W}_{\mathrm{in}}+\vect{b}_{\mathrm{in}}$ & $\shape{1,8,14336}$ \\
MLP activation $\mat{G}$ & $\phi(\mat{P})$ & $\shape{1,8,14336}$ \\
MLP update & $\mat{G}\mat{W}_{\mathrm{out}}+\vect{b}_{\mathrm{out}}$ & $\shape{1,8,4096}$ \\
block output & $\mat{U}+\Delta\mat{X}_{\mathrm{mlp}}$ & $\shape{1,8,4096}$ \\
final norm $\mat{H}$ & $\operatorname{Norm}_{f}(\mat{X}^{(L)})$ & $\shape{1,8,4096}$ \\
logits & $\mat{H}\mat{W}_{U}^{\mathsf T}+\vect{b}_{U}$ & $\shape{1,8,128256}$ \\
next-token logits & 마지막 position 선택 & $\shape{1,128256}$ \\
\bottomrule
\end{tabularx}
\end{table}

query, key와 value를 만들 때의 reshape는 model dimension을 head 수와 head dimension으로 나누고 attention 계산이 편한 축 순서 $\shape{B,H,T,d_h}$로 transpose하는 과정이다.
head output을 합칠 때는 이 축 변환을 되돌려 $\shape{B,T,Hd_h}$를 만들며, 이 예시에서는 $Hd_h=4096$이다.
GQA 모델에서는 key와 value의 head 수가 query head 수보다 작으므로 위 key/value shape가 달라진다.
이 차이는 \cref{chap:phase3}에서 정확히 전개한다.

\Needspace{38\baselineskip}
\section{Phase 2 학습 점검}

\begin{studybox}
\begin{enumerate}
  \item 길이 3, head dimension 2인 causal attention의 score, mask, softmax와 value 가중합을 손으로 계산한다.
  \item query의 position과 key/value의 position을 문장으로 설명하고 attention matrix의 행과 열을 표시한다.
  \item 한 head의 concatenation 표현을 head별 residual update의 합으로 바꿔 쓴다.
  \item attention이 읽는 $\mat{R}_{\mathrm{attn}}$과 MLP가 읽는 $\mat{R}_{\mathrm{mlp}}$가 왜 다른 tensor인지 설명한다.
  \item 작은 MLP의 input projection, nonlinearity와 output projection을 손으로 계산한다.
  \item MLP가 position을 직접 섞지 않으면서도 contextualized information을 처리할 수 있는 이유를 설명한다.
  \item pre-norm sequential block의 모든 중간 tensor shape를 적는다.
  \item position $t$의 logits와 학습 target $x_{t+1}$을 맞추고 loss mask가 필요한 위치를 표시한다.
  \item 마지막 residual update가 특정 token logit에 미치는 direct contribution을 unembedding vector와의 내적으로 계산한다.
  \item 같은 toy logits에 greedy, temperature, top-$k$와 top-$p$를 적용하고 남는 후보와 확률을 비교한다.
  \item 한 generation step에서 선택한 token이 다음 step의 조건부분포를 어떻게 바꾸는지 설명한다.
\end{enumerate}
\end{studybox}

\begin{limitbox}
이 장의 수식은 architecture, parameter와 실행 조건을 고정했을 때 각 tensor가 어떻게 계산되는지 명시한다.
그러나 특정 head가 왜 Obama와 출생지의 관계를 처리하는지, MLP activation이 어떤 feature를 나타내는지, 어느 경로가 Honolulu logit에 인과적으로 필요했는지는 수식만으로 정해지지 않는다.
높은 activation이나 attention weight는 분석할 후보를 제시할 수 있지만 causal importance를 증명하지는 않으며, 이런 주장은 intervention과 별도의 검증이 필요하다.
\end{limitbox}
